% =====================================================================
%  CAPÍTULO 1 — INTRODUÇÃO
% =====================================================================

\chapter{INTRODUÇÃO}
\label{cap:introducao}

O ciclismo virtual (\textit{indoor cycling}) transitou de uma simples
atividade alternativa para dias chuvosos para uma modalidade esportiva e
científica de alta relevância global. Atualmente, é reconhecida pela União
Ciclística Internacional (UCI) como uma modalidade oficial de
\textit{e-sports}, contando com campeonatos mundiais e atletas profissionais
dedicados \cite{coggan2010}. Este fenômeno foi impulsionado pelo surgimento de
rolos de treino inteligentes (\textit{smart trainers}), capazes de estabelecer
comunicação bidirecional com computadores e dispositivos móveis para simular a
resistência do terreno real em tempo real.

No entanto, a situação atual do ciclismo virtual é marcada por uma forte
centralização de mercado e barreiras econômicas severas. Plataformas líderes
como Zwift, Rouvy e TrainerRoad operam sob modelos de assinatura mensal
proprietários com taxas cobradas em dólar comercial, em torno de
\$20~USD/mês. Para atletas em países em desenvolvimento, como o Brasil, o
custo acumulado dessas assinaturas ultrapassa R\$~1.200,00 anuais (em média
\$240~USD/ano no plano mensal), o que constitui um fator de exclusão para
ciclistas amadores e atletas de base.

Além da barreira financeira, essas plataformas comerciais exigem hardware de
alto desempenho com placas gráficas dedicadas, inviabilizando o uso por
atletas que possuem apenas computadores de entrada ou dispositivos móveis
comuns. Adicionalmente, tais soluções dependem de conexão contínua com a
internet para computar métricas, o que impede treinos em ambientes
\textit{offline} ou com conectividade instável.

O \argus{} propõe-se a ser uma alternativa nacional, multiplataforma,
\textit{offline-first} e de código aberto, unindo simulação física interativa,
telemetria de precisão, assistência de inteligência artificial generativa
local e rastreamento de desgaste de componentes, sem taxas recorrentes e com
baixíssimos requisitos de hardware. Este documento constitui o registro do
software desenvolvido, apresentando sua arquitetura, implementação, resultados
e a evolução do projeto desde a versão registrada junto ao INPI.

\section{Contextualização: a situação atual do ciclismo virtual}
\label{sec:contextualizacao}

A seção a seguir apresenta o estado da arte do ciclismo virtual, comparando o
\argus{} com as principais plataformas do mercado e iniciativas de código
aberto, para situar o leitor quanto ao contexto tecnológico e mercadológico no
qual o software se insere.

\subsection{Estado da arte e referências a outros projetos}
\label{sec:estado-arte}

Para contextualizar a relevância e o posicionamento do \argus{}, apresenta-se
a seguir um comparativo com as principais plataformas do mercado e iniciativas
de código aberto:

\begin{itemize}
    \item \textbf{Zwift}: o líder absoluto do mercado foca em um ambiente
        virtual tridimensional massivo \textit{online}, com forte apelo social
        e gamificação. Suas limitações residem na ausência de modo
        \textit{offline}, custo recorrente elevado (\$19,99~USD/mês) e exigência
        de hardware potente para processamento gráfico 3D.
    \item \textbf{Rouvy}: diferencia-se por simular rotas reais por meio de
        vídeos gravados de estradas ao redor do mundo. Apresenta alta demanda
        de largura de banda de rede para transmissão de vídeo e sofre dos
        mesmos problemas de custo de assinatura fechada (\$19,99~USD/mês).
    \item \textbf{TrainerRoad}: foca exclusivamente em treinos estruturados
        baseados em faixas de potência (telemetria pura), sem simulação
        gráfica complexa. Suas métricas e planos de treino são proprietários e
        bloqueados por assinaturas de custo elevado (\$19,95~USD/mês).
    \item \textbf{GoldenCheetah}: software livre sob licença GPLv3 voltado à
        análise científica pós-treino. Embora possua ferramentas matemáticas
        excelentes (como o gráfico PMC e modelos de fadiga), sua interface é
        extremamente complexa, com uma curva de aprendizado íngreme, e ele não
        foi projetado para atuar como um simulador de rotas interativo e
        responsivo em tempo real com controle de rolos inteligentes.
    \item \textbf{Argus Cyclist}: posiciona-se no estado da arte ao preencher
        a lacuna entre a análise de dados científica (típica do GoldenCheetah)
        e a simulação de rotas interativa e gamificada (típica do Zwift),
        operando de forma totalmente gratuita, multiplataforma,
        \textit{offline-first} e com processamento local eficiente.
        Diferencia-se ainda por oferecer um treinador virtual baseado em IA
        generativa local (Ollama) \cite{ollama2024} para elaboração de treinos
        personalizados em linguagem natural e um sistema de rastreamento de
        desgaste de componentes que calcula automaticamente a vida útil
        restante de cada peça da bicicleta com base na distância percorrida.
\end{itemize}

\begin{table}[h]
    \centering
    \caption{Comparativo entre plataformas de ciclismo virtual}
    \label{tab:comparativo}
    \small
    \begin{tabular}{
        >{\raggedright\arraybackslash}p{2.8cm}
        >{\raggedright\arraybackslash}p{2.4cm}
        >{\raggedright\arraybackslash}p{2.2cm}
        >{\raggedright\arraybackslash}p{2cm}
        >{\centering\arraybackslash}p{1.1cm}
        >{\raggedright\arraybackslash}p{3cm}
    }
        \toprule
        \textbf{Plataforma} & \textbf{Licença} & \textbf{Execução} &
        \textbf{Hardware} & \textbf{Offline} & \textbf{Foco Principal} \\
        \midrule
        Zwift & Proprietária (\$19,99/mês) & Nuvem & Alto & Não &
        Gamificação e Social \\
        Rouvy & Proprietária (\$19,99/mês) & Nuvem & Médio-Alto & Não &
        Realismo de rotas reais \\
        TrainerRoad & Proprietária (\$19,95/mês) & Nuvem & Baixo & Não &
        Treinos Estruturados \\
        GoldenCheetah & Livre (GPLv3) & Local & Baixo & Sim &
        Análise Científica pós-treino \\
        \textbf{Argus Cyclist} & \textbf{Livre (GPLv3)} &
        \textbf{Local (Edge)} & \textbf{Baixo} & \textbf{Sim} &
        \textbf{Física + Telemetria + IA} \\
        \bottomrule
    \end{tabular}
\end{table}

\section{Problemática}
\label{sec:problematica}

O desenvolvimento do projeto \argus{} é motivado por quatro problemas centrais
identificados no ecossistema atual do ciclismo virtual:

\begin{enumerate}
    \item \textbf{Exclusão socioeconômica e barreiras de acesso}: o modelo de
        monetização em moeda estrangeira das plataformas de \textit{e-sports}
        de ciclismo inviabiliza o treinamento científico e estruturado para
        atletas de base e ciclistas amadores de menor poder aquisitivo no
        Brasil.

    \item \textbf{Dependência crônica de conectividade (\textit{cloud lock-in})}:
        a arquitetura das soluções comerciais centraliza a lógica do simulador
        e o processamento de métricas em servidores na nuvem. Quedas temporárias
        de conexão resultam no encerramento abrupto das atividades ou na perda
        de treinos inteiros, além de inviabilizar o uso em locais sem sinal de
        internet.

    \item \textbf{Vulnerabilidade de privacidade de dados}: a coleta de dados
        biométricos (frequência cardíaca contínua) e geográficos (coordenadas
        GPS de rotas importadas) é armazenada de forma compulsória em
        servidores privados estrangeiros, infringindo princípios de soberania
        de dados e privacidade por padrão (\textit{Privacy-by-Design}).

    \item \textbf{Monetização fisiológica abusiva}: métricas avançadas de
        controle de carga de treinos fundamentais para evitar o
        \textit{overtraining} e prescrever a recuperação (como PMC, TSS e CTL)
        dependem apenas de equações matemáticas consolidadas na literatura
        científica de fisiologia do esporte. Ainda assim, são vendidas como
        recursos exclusivos em softwares proprietários pagos.
\end{enumerate}

\section{Justificativa}
\label{sec:justificativa}

O projeto \argus{} justifica-se pela sua contribuição tecnológica, científica
e de impacto social no cenário esportivo e de engenharia de software.

\subsection{Relevância tecnológica e inovação arquitetural}
\label{sec:relevancia-tecnologica}

Do ponto de vista da Engenharia de Software, o projeto aborda o desafio de
criar um sistema multiplataforma de tempo real de baixíssima latência. A
escolha do Go (Golang) no \textit{backend} deve-se à sua capacidade nativa de
concorrência baseada em \textit{goroutines} e \textit{channels}, permitindo
realizar escaneamento e manter conexões estáveis com múltiplos sensores
Bluetooth Low Energy (BLE) utilizando a biblioteca TinyGo Bluetooth
\cite{tinygo2024} com perfis GATT específicos (frequência cardíaca
\cite{bluetooth2017}, potência e controle bidirecional FTMS) sem congelar a
interface de usuário.

A integração com o framework Wails \cite{wails2024} viabiliza um binário de
\textit{desktop} extremamente leve, utilizando a \textit{web engine} nativa do
sistema operacional para renderizar uma interface reativa em JavaScript
(Vite), registrando um consumo de memória inferior a \SI{200}{\mega\byte} ---
uma fração do demandado por soluções baseadas em Electron. Para dispositivos
móveis, a utilização do Capacitor.js \cite{capacitor2024} possibilita
reutilizar a mesma base de código do \textit{frontend} para compilar um
aplicativo nativo para Android, rodando um motor físico portado inteiramente
para JavaScript para simular e controlar os rolos diretamente de um celular
comum via antenas BLE integradas.

Inovadoramente, o sistema incorpora um módulo de IA generativa local que se
comunica com uma instância do Ollama \cite{ollama2024} rodando no próprio
dispositivo do usuário. Esta arquitetura elimina completamente a dependência de
serviços em nuvem para funcionalidades de \textit{coaching}, garantindo que
nenhum dado biométrico ou de treino seja transmitido a servidores externos. A
escolha de modelos pequenos (três bilhões de parâmetros, como o
\texttt{qwen2.5:3b}) viabiliza a execução mesmo em hardware modesto, e o tempo
de resposta estendido acomoda a inferência em CPU.

O sistema de controle de desgaste de componentes implementa um modelo de
depreciação linear simples, porém eficaz, onde cada componente cadastrado
(corrente, cassete, coroas, etc.) tem sua vida útil definida em quilômetros
pelo usuário, e o desgaste percentual é atualizado automaticamente ao final de
cada sessão de treino com base na distância percorrida. Este sistema segue uma
arquitetura em camadas bem definida com persistência em SQLite
\cite{sqlite2024} local, mantendo a coerência com o princípio
\textit{offline-first} do projeto.

\subsection{Relevância científica (modelagem física e fisiologia)}
\label{sec:relevancia-cientifica}

O sistema aplica modelagens matemáticas rigorosas de física e fisiologia do
esporte diretamente no dispositivo do usuário. O motor físico calcula a
aceleração e a velocidade virtual aplicando a conservação de energia sobre as
forças resistivas, conforme a equação de Martin \cite{martin1998}, detalhada
no Capítulo~\ref{cap:fundamentacao}. O módulo de fisiologia processa o vetor de
potência de \SI{1}{\hertz} para extrair métricas como Potência Normalizada
(NP), Fator de Intensidade (IF), \textit{Training Stress Score} (TSS),
\textit{Training Impulse} (TRIMP), além de compilar o
\textit{Performance Management Chart} (PMC) e o desvio cardiovascular
(\textit{Aerobic Decoupling}), fundamentadas na literatura de fisiologia do
esporte \cite{coggan2010}.

\subsection{Relevância social e inclusão esportiva}
\label{sec:relevancia-social}

Ao prover uma solução \textit{offline} de código aberto e totalmente gratuita
compatível com computadores modestos e celulares Android, o \argus{} atua como
um elemento de inclusão social. O projeto permite que ciclistas locais e
equipes amadoras de ciclismo utilizem equipamentos de rolo de treino
inteligente sem taxas mensais, democratizando o treinamento esportivo de alto
rendimento.

\section{Objetivos}
\label{sec:objetivos}

\subsection{Objetivo geral}
\label{sec:objetivo-geral}

Projetar, desenvolver e validar o \argus{}, um simulador multiplataforma de
ciclismo virtual de código aberto e gratuito, integrado a um sistema científico
de análise de telemetria esportiva \textit{offline-first}, garantindo
processamento em borda (\textit{edge}) de baixa latência e interoperabilidade
com hardware padrão do mercado.

\subsection{Objetivos específicos}
\label{sec:objetivos-especificos}

\begin{enumerate}
    \item \textbf{Conectividade de hardware}: implementar rotinas robustas em
        Go para escaneamento e conexão com periféricos via Bluetooth Low
        Energy (BLE), decodificando serviços GATT para cintas cardíacas
        (\texttt{0x180D}), sensores de potência (\texttt{0x1818}) e rolos de
        treino controlados (FTMS --- \texttt{0x1826}) \cite{bluetooth2017}.
    \item \textbf{Motor de simulação física}: desenvolver algoritmos que
        traduzam o gradiente das coordenadas de uma rota GPX
        \cite{gpx2004} em resistência de torque mecânico no rolo de treino
        inteligente, atualizando dinamicamente a resistência conforme o peso do
        ciclista e sua inclinação virtual.
    \item \textbf{Interface gráfica de alta performance}: construir um HUD
        (\textit{Heads-Up Display}) leve a 60~fps usando Canvas/WebGL,
        integrando o MapLibre GL JS \cite{maplibre2024} com renderização
        topográfica em 3D Terrain-RGB e rotas dinâmicas coloridas de acordo
        com a inclinação do terreno.
    \item \textbf{Análise fisiológica local}: implementar funções matemáticas
        para cálculo contínuo de métricas esportivas avançadas (NP, IF, TSS,
        TRIMP) e visualização do gráfico de evolução histórica de carreira PMC
        (\textit{Performance Management Chart}).
    \item \textbf{Assistente de treino por IA generativa local (IA Coach)}:
        integrar um assistente conversacional de treinamento esportivo
        utilizando modelos de linguagem de pequena escala (3B parâmetros)
        executados localmente via Ollama \cite{ollama2024}, capaz de
        interpretar o perfil do atleta e gerar planos de treino estruturados no
        formato ZWO (\textit{Zwift Workout}) com carregamento automático no
        modo ERG do simulador.
    \item \textbf{Controle de desgaste de componentes mecânicos}: implementar
        um sistema de cadastro e monitoramento de componentes da bicicleta com
        cálculo automático de desgaste percentual baseado na quilometragem
        acumulada, histórico de substituições com auditoria temporal e
        indicadores visuais codificados por cores (verde, amarelo, vermelho)
        para alertar o atleta sobre a necessidade de manutenção preventiva.
    \item \textbf{Persistência de dados confiável}: modelar um banco de dados
        local com SQLite \cite{sqlite2024} usando os padrões
        \textit{Repository} e \textit{Facade} para persistir usuários,
        atividades e recordes pessoais de forma 100\% \textit{offline} e
        estruturada.
    \item \textbf{Interoperabilidade de dados}: desenvolver rotinas para
        compilação e exportação de dados na estrutura de arquivo binário
        padrão \texttt{.FIT} (Garmin/ANT+ SDK) \cite{garmin2024} e integração
        de upload assíncrono via API do Strava \cite{strava2025} e do
        Intervals.icu \cite{intervals2024}.
    \item \textbf{Portabilidade multiplataforma}: compilar a aplicação para os
        sistemas Android, Windows, macOS e Linux utilizando tecnologias
        multiplataforma, adaptando o motor físico para JavaScript a fim de
        garantir operação autônoma, compatibilidade entre diferentes ambientes
        e baixa latência em dispositivos móveis e \textit{desktops}.
    \item \textbf{Gamificação saudável}: integrar um sistema de níveis e
        conquistas cardiovasculares persistentes com uma \textit{Trophy Room}
        interativa 3D/Canvas para engajamento dos atletas amadores.
\end{enumerate}

\section{Metodologia}
\label{sec:metodologia}

O trabalho adota uma metodologia de pesquisa aplicada, orientada ao
desenvolvimento tecnológico experimental, utilizando abordagens qualitativas e
quantitativas para validação dos resultados. O processo de engenharia de
software baseia-se em um modelo de desenvolvimento iterativo e incremental
(práticas ágeis), dividido em quatro fases essenciais.

\subsection{Fases do desenvolvimento}
\label{sec:fases}

\begin{enumerate}
    \item \textbf{Fase de levantamento e especificação}: estudo aprofundado dos
        perfis GATT Bluetooth SIG (\textit{Heart Rate Service},
        \textit{Cycling Power Service} e \textit{Fitness Machine Service}) e
        decodificação hexadecimal de \textit{frames} de telemetria. Análise da
        documentação do Garmin FIT SDK \cite{garmin2024} para modelar o
        codificador de arquivos binários e das APIs REST do Strava sob
        protocolo de segurança OAuth 2.0 \cite{strava2025}.

    \item \textbf{Fase de definição arquitetural}: desenho do padrão estrutural
        do sistema sob conceitos de \textit{Clean Architecture}. Separação
        rígida de responsabilidades em camadas (Apresentação, Casos de Uso,
        Serviços de Domínio e Persistência de Dados). Seleção da \textit{stack}
        multiplataforma ideal: \textit{backend} em Go (Golang) com
        \textit{bindings} Wails (Desktop) e empacotamento móvel nativo em
        Android via Capacitor.js \cite{capacitor2024}.

    \item \textbf{Fase de implementação e refatoração}: implementação dos
        serviços em Go: módulo de telemetria (\code{internal/service/ble}) com
        \textit{goroutines} concorrentes de escuta; módulo físico
        (\code{internal/service/sim}) com suporte a GPX e testes de FTP;
        persistência em banco local (\code{internal/repository/sqlite}) sob
        padrão \textit{Facade} (\code{usecase/storage\_facade}); e exportador
        esportivo (\code{internal/service/fit}). No \textit{frontend}
        JavaScript, implementa-se o \code{TelemetryEventBus} para desacoplar a
        atualização de tela da chegada de dados BLE, garantindo renderização
        fluida a 60~fps.

    \item \textbf{Fase de testes e validação experimental}: execução de testes
        funcionais e de integração utilizando equipamentos reais de ciclismo,
        conectando o \argus{} a rolos inteligentes (\textit{smart trainers}) e
        medidores de potência para validação prática do sistema em ambiente de
        uso real. A validação quantitativa consiste na medição da latência de
        telemetria (requisitada para operar em $<$ \SI{150}{\milli\second} do
        hardware até a tela) e em testes de conformidade estrutural,
        submetendo os arquivos binários \texttt{.fit} gerados diretamente ao
        Strava \cite{strava2025}, atestando a integridade estatística dos
        dados em relação ao padrão global da indústria.
\end{enumerate}

Dessa forma, a metodologia proposta garante a consistência científica no
desenvolvimento e assegura a qualidade física da simulação interativa,
resultando em um produto de software estável e utilizável por ciclistas.

\section{Organização do documento}
\label{sec:organizacao}

Este documento está organizado da seguinte forma: o Capítulo~\ref{cap:introducao}
apresenta a contextualização, problemática, justificativa, objetivos e
metodologia; o Capítulo~\ref{cap:fundamentacao} apresenta a fundamentação
teórica necessária à compreensão do sistema; o Capítulo~\ref{cap:registro-evolucao}
contextualiza o registro do software junto ao INPI e documenta a evolução do
projeto desde a versão registrada; o Capítulo~\ref{cap:arquitetura} descreve a
arquitetura do sistema; o Capítulo~\ref{cap:implementacao} detalha a
implementação dos módulos e funcionalidades; o Capítulo~\ref{cap:resultados}
apresenta os resultados e a validação do software; e o Capítulo~\ref{cap:conclusao}
tece as considerações finais e trabalhos futuros.
